%========================================================
% D1 — Expectation Values in the Density-Operator Formalism
%========================================================

\section{Expectation Values in the Density-Operator Formalism}

This derivation shows how the density-operator expression for expectation values generalizes the pure-state expression.

\medskip

For a pure state \( |\psi\rangle \), the expectation value of an observable \( O \) is

\begin{equation}
\langle O \rangle
=
\langle \psi | O | \psi \rangle.
\label{eq:pure_state_expectation_derivation}
\end{equation}

The associated density operator is

\begin{equation}
\rho
=
|\psi\rangle \langle \psi |.
\label{eq:pure_state_density_operator_derivation}
\end{equation}

Using this definition, the density-operator expression becomes

\begin{equation}
\langle O \rangle
=
\mathrm{Tr}(\rho O)
=
\mathrm{Tr}
\left(
|\psi\rangle \langle \psi| O
\right).
\label{eq:density_expectation_trace_form}
\end{equation}

To show that this reproduces the pure-state expression, expand the trace in an arbitrary orthonormal basis \( \{|k\rangle\} \):

\begin{align}
\mathrm{Tr}(\rho O)
&=
\sum_k
\langle k |
\psi\rangle
\langle \psi |
O
| k \rangle
\nonumber
\\
&=
\sum_k
\langle \psi |
O
| k \rangle
\langle k |
\psi\rangle
\nonumber
\\
&=
\langle \psi |
O
\left(
\sum_k |k\rangle \langle k|
\right)
|\psi\rangle
\nonumber
\\
&=
\langle \psi | O | \psi \rangle.
\label{eq:density_expectation_pure_equivalence}
\end{align}

Therefore,

\begin{equation}
\boxed{
\langle O \rangle
=
\mathrm{Tr}(\rho O)
}
\label{eq:density_expectation_general_result}
\end{equation}

reproduces the standard pure-state expectation value whenever \( \rho = |\psi\rangle \langle \psi| \).

\medskip

The advantage of the density-operator formalism becomes evident when the system is not described by a single pure state, but by a statistical ensemble.

Suppose the system is prepared in one of the pure states \( |\psi_k\rangle \) with probability \( p_k \), where

\[
p_k \geq 0,
\qquad
\sum_k p_k = 1.
\]

The density operator of the ensemble is

\begin{equation}
\rho
=
\sum_k
p_k
|\psi_k\rangle \langle \psi_k|.
\label{eq:ensemble_density_operator_derivation}
\end{equation}

The expectation value of \( O \) is then the statistical average of the corresponding pure-state expectation values:

\begin{align}
\langle O \rangle
&=
\sum_k
p_k
\langle \psi_k | O | \psi_k \rangle
\nonumber
\\
&=
\sum_k
p_k
\mathrm{Tr}
\left(
|\psi_k\rangle \langle \psi_k| O
\right)
\nonumber
\\
&=
\mathrm{Tr}
\left[
\left(
\sum_k
p_k
|\psi_k\rangle \langle \psi_k|
\right)
O
\right]
\nonumber
\\
&=
\mathrm{Tr}(\rho O).
\label{eq:ensemble_expectation_density_form}
\end{align}

Thus, the same expression

\begin{equation}
\langle O \rangle
=
\mathrm{Tr}(\rho O)
\label{eq:density_expectation_final}
\end{equation}

applies both to pure states and to statistical mixtures.

\medskip

This result is central for the simulator because it allows observable quantities to be evaluated through a unified mathematical interface, independently of whether the quantum state is represented as a pure state vector or as a density operator.

In particular, correlation observables of the form \( \sigma_i \otimes \sigma_j \), local polarization observables, and state diagnostics can all be computed through the same trace expression once the corresponding operator has been constructed.

\medskip

The density-operator expectation value therefore provides the mathematical bridge between deterministic pure-state evolution, ensemble-averaged descriptions, and effective non-unitary models arising from decoherence or statistical uncertainty.